\section{Закон инерции квадратичных форм}

\begin{definition}
Квадратичная форма $f(x, x)$ над полем $\mathbb{R}$ приводится к \textbf{нормальному виду}, если существует базис $E$, в котором она записывается как разность сумм квадратов:
\[
f(x, x) = \xi_1^2 + \dots + \xi_p^2 - \xi_{p+1}^2 - \dots - \xi_{p+q}^2,
\]
где $\xi_i$ --- координаты вектора $x$ в базисе $E$. Такой базис называется \textbf{нормальным базисом}.
\end{definition}

\begin{definition}
Числа $p$ и $q$ в нормальном виде квадратичной формы называются соответственно \textbf{положительным} и \textbf{отрицательным индексами инерции}. 
Число $n - (p + q)$ (где $n = \dim V$) называется \textbf{дефектом} (или нуль-индексом) квадратичной формы. Сумма $p + q$ равна рангу матрицы квадратичной формы.
\end{definition}

\begin{theorem}[Закон инерции]
Для любой квадратичной формы над полем $\mathbb{R}$ её положительный индекс инерции $p$, отрицательный индекс инерции $q$ и дефект $n - (p + q)$ являются инвариантами и не зависят от выбора нормального базиса.
\end{theorem}

\begin{proof}
Докажем инвариантность $p$ и $q$ (откуда автоматически последует инвариантность дефекта, так как $n$ и ранг $p+q$ инвариантны).

Предположим противное: пусть существуют два нормальных базиса $E = \{e_1, \dots, e_n\}$ и $E' = \{e'_1, \dots, e'_n\}$, в которых форма имеет индексы $(p, q)$ и $(r, s)$ соответственно, причем $p > r$.

Рассмотрим два подпространства:
\begin{itemize}
    \item $L_1 = \mathcal{L}(e_1, \dots, e_p)$ --- подпространство, натянутое на первые $p$ векторов базиса $E$. Его размерность $\dim L_1 = p$.
    \item $L_2 = \mathcal{L}(e'_{r+1}, \dots, e'_n)$ --- подпространство, натянутое на последние $n - r$ векторов базиса $E'$. Его размерность $\dim L_2 = n - r$.
\end{itemize}

Оценим размерность суммы этих подпространств:
\[
\dim(L_1 + L_2) = \dim L_1 + \dim L_2 - \dim(L_1 \cap L_2) \le n.
\]
Отсюда получаем оценку для размерности пересечения:
\[
\dim(L_1 \cap L_2) \ge p + (n - r) - n = p - r.
\]
Так как по предположению $p > r$, то $p - r > 0$, следовательно, $\dim(L_1 \cap L_2) > 0$. Это означает, что существует ненулевой вектор $x_0 \in L_1 \cap L_2$.

Рассмотрим значение квадратичной формы на этом векторе $x_0$ в двух базисах:
\begin{enumerate}
    \item Так как $x_0 \in L_1$, его координаты $\xi_{p+1}, \dots, \xi_n$ в базисе $E$ равны нулю. Поскольку $x_0 \neq 0$, хотя бы одна из координат $\xi_1, \dots, \xi_p$ отлична от нуля. Следовательно:
    \[
    f(x_0, x_0) = \xi_1^2 + \dots + \xi_p^2 > 0.
    \]
    \item Так как $x_0 \in L_2$, его координаты $\eta_1, \dots, \eta_r$ в базисе $E'$ равны нулю. Следовательно:
    \[
    f(x_0, x_0) = -\eta_{r+1}^2 - \dots - \eta_{r+s}^2 \le 0.
    \]
\end{enumerate}
Мы пришли к противоречию ($f(x_0, x_0) > 0$ и $f(x_0, x_0) \le 0$ одновременно). Значит, наше предположение $p > r$ неверно, и должно выполняться $p \le r$. 

В силу симметрии рассуждений (поменяв местами $E$ и $E'$), получаем $r \le p$. Следовательно, $p = r$.
Поскольку ранг формы $p + q = r + s$ инвариантен, из $p = r$ немедленно следует $q = s$.
\hfill$\blacksquare$
\end{proof}
